\section{Hidden cooling advantages approaching unity with dimension}
\label{sec:highdimamplitude}
The sharp four-state bound does not give a dimension-independent bound.
Cloning the high-energy state produces an increasingly large normalization
change along the initial Gibbs curve, while retaining a uniformly isolated
slow relaxation mode.

\begin{theorem}[A dimension-dependent achievable advantage]
\label{thm:highdimamplitude}
For every integer $n\geq4$ there is a family of strictly positive,
reversible $n$-state generators at bath inverse temperature one, with
all energies distinct and all original $F_{ijk},G_{ijk}<0$, such that
finite $0<\beta_h<\beta_c<1$ satisfy
\[
1-\frac{|b_2(\beta_h)|}{|b_2(\beta_c)|}
\longrightarrow 1-\frac3n.
\]
The slow eigenvalue tends to $-3$ and its separation from the next mode
is greater than $1/4$ for all sufficiently small family parameters.
The corresponding leading asymptotic KL advantage tends to $1-9/n^2$.
Consequently no dimension-independent upper bound less than one can
hold for these amplitude or leading KL advantages under the full
triplet premise. Sharpness is asserted only for $n=4$, by the preceding
four-state theorem; no matching upper bound is claimed for $n>4$.
\end{theorem}
\begin{proof}
Fix $n$, put $m=n-3$, and choose strictly decreasing clone weights
\[
r_j=\frac{2^{m-j}}{2^m-1},\qquad j=1,\ldots,m,\qquad \sum_jr_j=1.
\]
Start with the quarter-limit family $M_\varepsilon$, its normalized bath
weights $q_i$, and its left slow vector $w$ from
Theorem~\ref{thm:quarterlimit}. Replace its fourth state by $m$ states
$c_1,\ldots,c_m$ with weights $q_4r_j$. Equivalently the unnormalized
bath weights are
\[
\widehat q^{(n)}=(1,1-\varepsilon/3,3/5,
                  \varepsilon^2r_1,\ldots,\varepsilon^2r_m).
\]
Their sum remains the four-state partition function. For small positive
$\varepsilon$ they give $n$ distinct ordered energies
$E_i=-\log\widehat q_i^{(n)}$.

Keep the original rates among the first three states. From each clone to
low state $i\leq3$ set $a_{i,c_j}=a_{i4}$, and in the reverse direction
set $a_{c_j,i}=r_j a_{4i}$. Within the clone block use conditional reset
at rate four: $a_{c_j,c_k}=4r_j$ for $j\ne k$.
Set all diagonals by conservation, giving a generator $\overline M$.
All rates are positive and detailed balance is immediate.
Lumping the clone probabilities exactly recovers $M_\varepsilon$.
The additional $(m-1)$-dimensional subspace consists of probability
perturbations supported on the clones whose sum is zero. On this
subspace the generator is scalar, with eigenvalue
\[
-4-\sum_{i=1}^3a_{i4}<-4.
\]
Thus $\overline M$ retains the simple slow eigenvalue $-3$, gap exceeding
$1/3$, and left eigenvector
\[
\overline w=(1,-4/5,w_3,-\varepsilon,\ldots,-\varepsilon).
\]
It has zero bath mean and weighted norm at most one for small
$\varepsilon$.

Here are all the original-rate triplet checks before a small repair.
Triplets in the first three states are unchanged. For two low states
$i<j\leq3$ and a clone $c_l$,
\[
G_{ijc_l}=G_{ij4},\qquad
F_{ijc_l}=F_{ij4}+
 \frac{(r_l-1)q_4}{2q_j}G_{ij4}.
\]
For one low state $i$ and two clones $c_j<c_k$,
\[
G_{ic_jc_k}=r_j(a_{4i}-4),\qquad
F_{ic_jc_k}=(r_j+r_k/2)(a_{4i}-4).
\]
All-clone triplets have $F=G=0$.
The original rates satisfy $0<a_{i4}<2$ throughout
$0<\varepsilon\leq1/30$. This bound is proved by three exact Bernstein
checks on the rational generator. Their finite limits are
\[
(a_{14},a_{24},a_{34})\longrightarrow(220,220,132)/169.
\]
Consequently $a_{4i}=a_{i4}q_4/q_i\leq(10/3)\varepsilon^2<4$,
so the one-low triplets are strictly negative. For the two-low/clone
triplets, $G_{ij4}\leq0$ and $-G_{ij4}\leq a_{j4}<2$, while
$q_4/q_j\leq(5/3)\varepsilon^2$. Thus the possible positive defect in
$F$ is at most $(5/3)\varepsilon^2$. Every $G$ is already nonpositive.

Define an auxiliary generator $H$ at the new bath weights. On the first
three states its reduced conductances are
$k_{12}=100,k_{13}=10,k_{23}=11$. Between low state $i$ and any clone
use $k_{i,c_j}=i$; among ordered clones use constant downward rate four,
with reverse rates fixed by detailed balance. All its tests are strictly
negative throughout $0<\varepsilon\leq1/30$.
The first-three tests are exactly those of the certified auxiliary
generator in the quarter-limit construction. Write
$Q=13/5-\varepsilon/3+\varepsilon^2$ for the unchanged partition sum.
For low pairs $12,13,23$ and clone $c_l$, the $F$ values of $H$ are
\[
\frac{-1-98\varepsilon/3-\varepsilon^2r_l/2}{Q},\qquad
\frac{-24/5-\varepsilon^2r_l}{Q},\qquad
\frac{-21/5+3\varepsilon-\varepsilon^2r_l/2}{Q}.
\]
Each is at most $-1/3$ on the interval, and the $G$ values are
strictly negative. For one low state $i$ and clones $j<k$ the tests are
\[
G=i q_{c_j}-4<0,\qquad F=G(1+q_{c_k}/(2q_{c_j}))<0.
\]
For three clones they are $G=4(t-1)<0$ and
$F=4(t-1)(1+u/2)<0$, since the $r_j$ are distinct and ordered.

Now put $\widetilde M=\overline M+6\varepsilon^2H$.
Linearity makes every $G$ strict. In two-low/clone triplets,
\[
F(\widetilde M)\leq(5/3)\varepsilon^2
                 -6\varepsilon^2/3=-\varepsilon^2/3<0.
\]
All other $F$ tests are already negative or acquire a strict negative
contribution from $H$. Thus all tests are strict on the entire
interval $0<\varepsilon\leq1/30$.

The repair does not spoil the spectrum or limiting overlaps. Let
$D_n=\operatorname{diag}(q^{(n)})$.
The symmetric matrix $D_n^{-1/2}HD_n^{1/2}$ has operator norm at most
\[
C_n=2\{103+4(m-1)\}.
\]
Indeed all rates leaving a low state sum to at most $103$, and those
leaving a clone sum to at most $3+4(m-1)$; each upward internal clone
rate is at most its downward rate four. Gershgorin applied to the
column generator and reversibility place its eigenvalues in
$[-C_n,0]$. This argument avoids any assumption of uniformly bounded
reduced conductances inside the increasingly rare clone block.
Weyl's inequality gives slow separation at least
$1/3-12C_n\varepsilon^2>1/4$ for sufficiently small $\varepsilon$.
The isolated eigenvector perturbation estimate in the quarter-limit
proof gives, with consistent orientation and equal weighted norm,
\[
\|\widetilde w-\overline w\|_{q^{(n)}}=O_n(\varepsilon^2).
\]
Since $q_i^{(n)}\geq r_m\varepsilon^2/3$ for small $\varepsilon$,
\[
\max_i|\widetilde w_i-\overline w_i|=O_n(\varepsilon).
\]
Every thermal overlap changes by at most this quantity, uniformly in
initial inverse temperature. The constants may depend on the fixed
integer $n$; no joint limit with $n$ and $\varepsilon$ is used.

Finally put $L=\log(1/\varepsilon)$ and choose
$\beta_h=L^{-2}$, $\beta_c=L^{-1/2}$. At $\beta_h$ every unnormalized
Gibbs weight tends to one. At $\beta_c$ each clone weight
$\exp(-2\sqrt L)r_j^{1/\sqrt L}$ tends to zero, while the first three
weights tend to one. Since the first three limiting slow components
sum to $-2/15$ and all clone components tend to zero,
\[
b_h\longrightarrow-\frac{2}{15n},\qquad
b_c\longrightarrow-\frac2{45}.
\]
The repair error tends to zero, so the ratio tends to $3/n$.
Both limiting overlaps are nonzero for each fixed $n$, and $3/n<1$,
which proves an admissible strict cooling witness for sufficiently
small positive $\varepsilon$. The leading KL ratio is the square of
the amplitude ratio, giving the stated KL advantage. Letting $n$
grow after taking these fixed-dimensional limits proves the absence
of a dimension-independent bound below one.
\end{proof}

The checker \path{research/high_dimension_amplitude_certificate.py}
records the exact cloning identities and six Bernstein sign certificates.
It also constructs the rational five-state member with
$\varepsilon=1/1000$ and clone weights $(2/3,1/3)$, verifies every strict
triplet test, positivity, conservation and detailed balance, and certifies
a perturbed slow gap of at least $124037/375000>1/4$.
This illustrative finite member verifies the generator construction; it
is not presented as a separately interval-certified thermal witness.
The all-dimensional thermal advantage is established by the limit proof.

There is also a separately certified finite five-state thermal witness.
For $m=2$, every outgoing rate sum of $H$ is at most $100$, so its
symmetric operator norm is at most $200$. Write
$\eta=1200\varepsilon^2$ for the perturbation bound. For
$\varepsilon\leq1/200$, the gap stays above
$1/3-2\eta\geq41/150>1/4$, and the oriented eigenvector estimate gives
\[
\|\widetilde w-\overline w\|_{q^{(5)}}
\leq\frac{\sqrt2\,\eta}{1/3-\eta}<6000\varepsilon^2.
\]
The normalization is the same weighted norm as $\overline w$, which is
at most one. Since the smallest normalized bath weight is at least
$\varepsilon^2/9$, every component and thermal overlap changes by less
than $18000\varepsilon$.
At $\varepsilon=10^{-50}$, $\beta_h=1/10000$ and $\beta_c=1/50$,
integer-power root bounds, combined with that error, enclose the repaired
amplitude advantage between $37.6778168708\%$ and $37.6778168731\%$.
The exact rational endpoints and reproducible computation are in
\path{research/five_state_large_advantage_certificate.py} and its JSON
receipt. This witness exceeds the four-state supremum. Its very large
energy span remains a mathematical construction, not evidence of a
practical cooling protocol.
